%!TeX root=../carlo_master.tex
\chapter{Introduction}\label{sec:intro}

In supervised learning, high-capacity models such as Deep Neural Networks and Gradient Boosted Machines are trained to approximate the conditional distribution $\mathbb{P}(Y \mid X)$ of a target variable $Y$ given input features $X$. In practice, however, the goal is rarely the probability estimate itself, but a discrete decision $\hat{y}$ that triggers a specific action.

Many applied settings carry asymmetric costs. In medical diagnostics, fraud detection, and industrial fault monitoring, a false negative often incurs a penalty orders of magnitude greater than a false positive. The default $\arg\max$ rule, which predicts the class with the highest estimated probability, implicitly minimizes the $0/1$ loss and is suboptimal under such asymmetry. It is likewise suboptimal when the operational objective is a non-decomposable performance metric such as $F_\beta$ or the Matthews correlation coefficient, which does not reduce to a per-instance expected cost.

Bayesian decision theory prescribes the answer: if the estimated probabilities are calibrated, the minimum-risk action is the class minimizing the expected conditional risk under the operational cost. The qualifier matters. Calibration and discriminative power are separate properties of an estimator, and a classifier can rank instances correctly by conditional class probability while emitting miscalibrated probability values. Modern classifiers routinely do so, with overconfident scores from deep networks at one end and conservative estimates from highly regularized models at the other. Applying the analytical Bayes rule to miscalibrated estimates yields excess risk over the Bayes optimum.

This dissertation addresses that gap by separating the \textit{inference stage} (probability estimation) from the \textit{decision stage} (risk minimization), and treating the decision stage as a learnable component fitted directly on a held-out validation partition rather than read off the analytical plug-in formula. Concretely, the binary threshold or the multiclass simplex partition is tuned by empirical risk minimization on validation data. The decision object is small (a scalar in binary, a low-parameter simplex partition in multiclass, with dimension polynomial in $K$) compared with the inference estimator that supplies its inputs, so the optimization is in principle learnable from a modest label budget.

\paragraph{Contributions} The main contributions of this work are:
\begin{enumerate}
    \item A unified treatment of multiclass partition-optimization strategies within a common partition-matrix framework (Chapter~\ref{sec:partition_optimization}), with a consolidated sample-complexity analysis of the parallel-shift, scaling, full-matrix, and Voronoi classes (Section~\ref{sec:multiclass_sample_complexity}).
    \item Voronoi tessellations as a parameterization of multiclass simplex partitions: under a dominant label constraint they exhibit a strict population risk separation from the parallel-shift class on a witness construction, and a sequential sampler preserves the constraint during search (Chapter~\ref{sec:voronoi_tessellation}).
    \item A large-scale empirical comparison of decision-side optimization strategies across fourteen binary and eleven multiclass datasets, ten classifier architectures, and ten binary and twelve multiclass objectives spanning cost matrices and non-decomposable maximization metrics (Chapters~\ref{sec:binary_experiments} and~\ref{sec:multiclass_experiments}).
    \item A characterization of how the relative excess risk of the analytical Bayes-optimal rule under estimation misspecification varies with the number of classes, supported by two synthetic perturbation modes and a dimensional-dilution candidate mechanism (Chapter~\ref{sec:misspec_sensitivity_K}).
    \item The open-source Python library \texttt{threshold\_opt} (\url{https://github.com/vitaliset/threshold-opt}, on PyPI as \texttt{threshold-opt}), which implements every decision strategy compared in this dissertation, the multiclass partition strategies included, with scikit-learn estimator conventions; the same repository holds the experiment code behind every reported result.
\end{enumerate}

\paragraph{Dissertation organization} The work is organized as follows. Chapter~\ref{sec:theory} establishes the decision-theoretic framework and defines the Bayes Optimal decision rule. Chapter~\ref{sec:binary} explores the binary case, contrasting analytical and empirical thresholding strategies. Chapter~\ref{sec:binary_experiments} presents the corresponding empirical results. Chapter~\ref{sec:multiclass} generalizes these concepts to multiclass simplex partitioning, and Chapter~\ref{sec:partition_optimization} surveys approaches to partition optimization. Chapter~\ref{sec:voronoi_tessellation} proposes Voronoi tessellations as a decision hypothesis class and analyzes their theoretical properties. Chapter~\ref{sec:multiclass_experiments} presents the multiclass empirical evaluation. Chapter~\ref{sec:misspec_sensitivity_K} probes the sensitivity of the analytical Bayes-optimal rule to estimation misspecification as a function of the number of classes, and proposes a candidate dimensional-dilution mechanism for the observed $K$-monotone flattening. Chapter~\ref{sec:conclusion} concludes with a discussion on the estimation-decision trade-off.
